% Propietario conservado: /Users/ruben/Documents/New project/output/REVISION_ARGUMENTAL_APERTURAS_CIERRES_20260915/VI/delivery/MANUSCRIPT_VI_ES_EN_COMPLETE/source_es/sections/08_familias_y_compuestos.tex
\section{Réalisations sectorielles, composition matérielle et classification spectrale}
\label{vi:vi:sec:familias-compuestos}

L'atlas de la section~\ref{vi:vi:sec:atlas} et le lecteur de la
section~\ref{vi:vi:sec:operador-masa} ont des fonctions complémentaires. Le
premier classe supports, profondeurs et orbites ; le second conserve la
chronologie de chaque trajectoire et publie son opérateur de masse. Leur application
à une famille exige de maintenir les deux structures. Une égalité entre
valeurs de masse n'identifie pas des routes ; une égalité entre cellules n'identifie pas non plus
des histoires, des représentations ou des états composites.

Cette section développe les réalisations sectorielles de cette
construction. Les constantes internes, la paire angulaire, l’échelle absolue
et les lecteurs bidirectionnels sont ceux déjà construits à partir
d’APP--TRIT--TPK. La représentation physique précise quelle information de
l’état est observée ; la comparaison métrologique s’effectue après avoir fixé
cette représentation. L’exposé conserve, en particulier, la différence
entre un opérateur de fibre, une évaluation exacte pour une signature déclarée
et une raie spectrale individualisée par un couplage.

\subsection{Descripteurs sectoriels et critère spectral d'individualisation}
\label{vi:vi:subsec:descriptor-familia}

Pour un état élémentaire, nous écrivons
\[
 X=(\mathfrak c,\rho,g,f,q,\epsilon,\varsigma,\mathfrak a).
\]
Ici \(\mathfrak c\) est la classe sectorielle, \(\rho\) la représentation,
\(g\) la génération physique, \(f\) la saveur, \(q\) la charge,
\(\epsilon\) l'orientation de feuillet, \(\varsigma\) la parité statistique
et \(\mathfrak a\) les données du couplage. La profondeur combinatoire
de l'atlas et la génération physique restent des coordonnées de types
distincts. La saveur distingue les composantes d'une représentation de
famille ; la couleur distingue les composantes de la représentation chromatique ;
le feuillet et son orientation conservent la provenance opératorielle de la
trajectoire. Aucune de ces étiquettes ne remplace le registre d'événements.

Soit \(\mathcal P_{\rm adm}\) l’ensemble des descripteurs physiques admissibles,
soit \(\mathcal R_{\rm adm}\) le registre des routes enrichies produit par
l’atlas et soit \(\mathcal D\) l’espace des descripteurs d’incidence qui
conservent classe sectorielle, représentation, profondeur, charge, couleur,
orientation et données de couplage. Nous définissons
\[
 d_{\mathcal P}:\mathcal P_{\rm adm}\longrightarrow\mathcal D,
 \qquad
 d_{\mathcal R}:\mathcal R_{\rm adm}\longrightarrow\mathcal D,
\]
où \(d_{\mathcal P}(X)\) publie les exigences structurelles de l'état
\(X\) et \(d_{\mathcal R}(\gamma)\) publie les mêmes coordonnées conservées
par la route \(\gamma\). La fibre, l'orbite ou la multisection produite par le
descripteur est le produit fibré
\begin{equation}
 F_X=d_{\mathcal R}^{-1}\!\bigl(d_{\mathcal P}(X)\bigr)
 =\mathcal R_{\rm adm}\times_{\mathcal D}\{d_{\mathcal P}(X)\}.
 \label{vi:vi:eq:fibra-descriptor}
\end{equation}
Pour les descripteurs réalisables de l’article, \(F_X\) est finie et non vide.
Sur
\(\mathcal H_X=\ell^2(F_X)\otimes V_{\rho_X}\) agissent les deux opérateurs
\(M_X^{\beta,D}\) et \(M_X^{\beta,\Omega}\) de la
section~\ref{vi:vi:sec:operador-masa}. Comme ils sont diagonaux dans la base de
routes, leur mesure spectrale conjointe est
\begin{equation}
 E_X(B)=\sum_{\gamma\in F_X}
 \mathbf1_B\bigl(m_D(\gamma),m_\Omega(\gamma)\bigr)
 |\gamma\rangle\langle\gamma|\otimes I_{V_{\rho_X}},
 \qquad B\subset\mathbb R_{>0}^2.
 \label{vi:vi:eq:familias-pvm}
\end{equation}
Conserver cette mesure permet de distinguer deux fibres même si leurs moyennes
coïncident. Une fonctionnelle scalaire des deux lecteurs ne s’applique que
lorsqu’elle fait partie de la représentation déclarée. Si le couplage
conserve plusieurs valeurs, la sortie est un spectre avec multiplicités.

\begin{proposition}[Critère de raie spectrale]
\label{vi:vi:prop:familias-linea}
Soit \(M\) un opérateur autoadjoint sur une fibre finie et \(P\) un
projecteur de préparation. Le canal publie exclusivement la valeur \(m\)
si et seulement si
\[
 E_M(\{m\})P=P,
 \qquad\text{de manière équivalente}\qquad (M-mI)P=0.
\]
L’égalité \(PMP=mP\), à elle seule, détermine une compression scalaire,
mais ne garantit pas que le support spectral préparé soit ponctuel.
\end{proposition}
\begin{proof}
Par le théorème spectral, \(\ker(M-mI)\) est l’image de
\(E_M(\{m\})\). Les deux premières conditions expriment que
\(\operatorname{im}P\) est contenue dans cet espace propre. Pour vérifier
la dernière distinction, prenons \(M=\operatorname{diag}(1,3)\) et le projecteur
sur \((1,1)/\sqrt2\). Alors \(PMP=2P\), tandis que la distribution
de masse préparée a pour support \(\{1,3\}\), non \(\{2\}\).
\end{proof}

L'individualisation conserve ainsi la portée de chaque résultat. La fibre
et l'opérateur sont des objets déterminés ; le passage à une droite utilise le
projecteur et la condition spectrale indiqués. Cette formulation s'appliquera
de la même manière aux familles chargées, neutres et composites.

\subsection{Leptons chargés : centre électronique et multisections de génération}
\label{vi:vi:subsec:leptones-cargados}

Les trois multisections chargées possèdent les rangs \(1,8,16\) aux profondeurs
\(G_0,G_2,G_4\). Elles partagent la représentation de charge et la trivialité
de couleur, mais ni le rang du support ni le registre de la trajectoire.
L’électron occupe le centre \((5,5)\) ; les réalisations muonique et
tauique conservent leurs multisections jusqu’au couplage
spectral. La conjugaison de leurs représentations est traitée au moyen des
opérateurs de la section~\ref{vi:vi:sec:conjugacion}, non au moyen d’un nouveau
choix de masses.

La coordonnée électronique basale \(\mathfrak e_0\) fournit la carte
angulaire intermédiaire
\begin{equation}
 \mathfrak m_X^{(0)}=\mathfrak e_0
 \exp\left(n_A(X)x+\frac{s_C(X)}6y\right),
 \qquad x=A_{\mathrm{rad}},\quad y=C^*_{\mathrm{rad}}.
 \label{vi:vi:eq:carta-angular-familias}
\end{equation}
Dans cette formule, \(\mathfrak e_0\) est la section basale et non la fermeture
électronique bêta. Le retour bidirectionnel appartient à l’étape ultérieure
de la construction. Les deux étapes conservent leur propre normalisation.

Les registres nominaux récupérés pour les deux familles non centrales
publient
\[
 (n_A,s_C)_\mu=(42,-3),\qquad (n_A,s_C)_\tau=(64,0),
\]
et leurs signatures complètes déclarées sont
\begin{equation}
 \sigma_\mu=(42,-3,8,0,2),\qquad
 \sigma_\tau=(64,0,25,-1,6).
 \label{vi:vi:eq:firmas-muon-tau}
\end{equation}
Pour expliquer l’effet de chaque coordonnée, notons \(\mathcal E_j\)
les lectures successives d’une signature déclarée :
\begin{align}
 \mathcal E_0(\sigma)&=n_Ax+s_Cy/6,\notag\\
 \mathcal E_1(\sigma)&=\mathcal E_0(\sigma)+k_{\Delta_4}\Delta_4,\notag\\
 \mathcal E_2(\sigma)&=\mathcal E_1(\sigma)+b_{120}s_{120},\notag\\
 \mathcal E_3(\sigma)&=\mathcal E_2(\sigma)+b_\zeta\zeta_{270},
 \qquad\zeta_{270}=135\Delta_4^2.
 \label{vi:vi:eq:etapas-e3-sectoriales}
\end{align}
Le symbole \(\mathcal E_3\) désigne ici un exposant d’évaluation,
non une espèce ni un mode supplémentaire. Le supplément bêta qui conserve ces
signatures emploie la normalisation explicite
\(m_X^{[3]}=m_e^\beta\exp\mathcal E_3(\sigma_X)\).
C’est une évaluation relative aux signatures du supplément ; elle ne remplace pas
la construction route par route de \(M_X^{\beta,D}\) et de
\(M_X^{\beta,\Omega}\).

La différence des signatures donne, sans approximation,
\begin{equation}
 \log\frac{m_\tau^{[3]}}{m_\mu^{[3]}}
 =22x+\frac12y+17\Delta_4-s_{120}+4\zeta_{270}.
 \label{vi:vi:eq:razon-tau-muon}
\end{equation}
La démonstration consiste à soustraire les exposants de
\eqref{vi:vi:eq:etapas-e3-sectoriales}. Le terme angulaire, la mémoire
torsionnelle et les deux lectures de retour restent identifiés dans le
résultat. Les évaluations conservées sont
\begin{align*}
 m_\mu^{[3]}&=105.658360294435829303\ldots\ \mathrm{MeV}/c^2,\\
 m_\tau^{[3]}&=1776.860917631432295595\ldots\ \mathrm{MeV}/c^2.
\end{align*}
Leur statut documentaire est celui d’une évaluation exacte lorsque la signature récupérée est fixée.
La sélection physique de cette signature au sein de la multisection et le critère
de raie spectrale sont des opérations différentes et sont conservées comme telles.

Au centre électronique, le rang est un et les lecteurs coïncident en
\((80,54,6)\). L'histoire active
\((2,2,5,-4,-3,-2)^2\) conserve plus d'information que son enveloppe
historique \((4,3,5,-4,-3,-5)^2\), bien que les deux partagent le motif de
signes. Cette distinction relie la particularité électronique à la
mémoire de position des événements : les familles ne se définissent pas seulement par
la forme visuelle d'un mot de signes.

\subsection{Neutrinos : neutralité, spectre de masse et transport de saveur}
\label{vi:vi:subsec:neutrinos-familias}

Le support neutrinique possède une dimension \(20=2+4+6+8\), répartie entre
quatre profondeurs. Ses trois étiquettes de saveur
\(\nu_e,\nu_\mu,\nu_\tau\) et ses trois valeurs propres possibles de masse
appartiennent à la réalisation physique du secteur, non au recensement des
profondeurs. La neutralité électrique conserve l’orientation, la
mémoire de feuillet et l’holonomie spinorielle. C’est pourquoi elle n’élimine pas les
variables qui interviennent dans le mélange.

Le registre neutre sépare les contributions de deux demi-cycles. Si
\(\mathcal I_m\) est le bloc de \(9\) itérations de la phase \(m\),
\(\theta_t\) marque la visite à \(9\), \(\epsilon_t\) conserve le feuillet
et \(\omega_t\) son orientation de couronne, le lecteur local est
\[
 T_\nu(m)=\sum_{t\in\mathcal I_m}\omega_t\epsilon_t\theta_t,
 \qquad e_m^\nu=T_\nu(m)-T_\nu(m+6).
\]
Cette différence maintient l’information orientée même lorsque sa somme
totale est nulle. L’opérateur complet utilise en outre les prédécesseurs
torsionnels, les vacances et le transport entre feuillets. Le passage au
sous-espace physique requiert le projecteur neutre de rang \(3\),
\[
 M_\nu=P_\nu\mathcal M_\nu P_\nu
 \quad\text{sur }\operatorname{im}P_\nu.
\]
Les valeurs propres de cette compression sont les masses dans l’échelle interne
action--horloge--masse. Les angles de mélange agissent entre leurs repères ;
ils ne remplacent pas l’échelle et n’engendrent pas rétrospectivement les valeurs propres.
La compression coïncide avec la restriction de l’opérateur complet lorsque
\(\operatorname{im}P_\nu\) est un sous-espace réducteur de
\(\mathcal M_\nu\).

Le registre permet de réunir un noyau de Gram avant de diagonaliser.
Pour chaque type de donnée
\(r\in\{\mathrm{car},\mathrm{hoj},\mathrm{ori},\mathrm{vac},\mathrm{hol}\}\),
soit \(\Xi_r(c)\) le vecteur produit par la cellule ou la route \(c\).
Nous définissons
\begin{equation}
 K_{\mathrm{reg}}(c,d)=
 \sum_{r:N_r>0}\frac{\langle\Xi_r(c),\Xi_r(d)\rangle}{N_r},
 \qquad N_r=\sum_a\|\Xi_r(a)\|^2.
 \label{vi:vi:eq:gram-neutrinos}
\end{equation}
L’omission d’un terme de norme totale nulle est fixée avant
l’évaluation. Elle conserve un domaine bien défini et n’introduit pas d’ajustement.

\begin{proposition}[Positivité et rang du registre neutre]
\label{vi:vi:prop:gram-neutrinos}
Le noyau \(K_{\mathrm{reg}}\) est hermitien et semi-défini positif.
Pour une agrégation tridimensionnelle \(G\), la condition \(\det G\ne0\)
certifie que les données agrégées séparent trois directions. Pour toute
matrice carrée inversible \(G\), son facteur polaire
\[
 U_G=G(G^*G)^{-1/2}
\]
est unitaire. Ce facteur ne détermine pas à lui seul les repères spectraux
des opérateurs chargé et neutre.
\end{proposition}
\begin{proof}
Pour des coefficients \(z_c\), la forme quadratique de
\eqref{vi:vi:eq:gram-neutrinos} est
\[
 \sum_{c,d}\overline z_cK_{\mathrm{reg}}(c,d)z_d
 =\sum_{r:N_r>0}N_r^{-1}
 \left\|\sum_c z_c\Xi_r(c)\right\|^2\ge0.
\]
Le rang est la dimension de l’image des vecteurs agrégés. Pour la
dernière affirmation,
\(U_G^*U_G=(G^*G)^{-1/2}G^*G(G^*G)^{-1/2}=I\) ;
l’inversibilité et l’égalité des dimensions donnent également \(U_GU_G^*=I\).
La formule utilise la matrice \(G\), tandis qu’un repère spectral exige de
diagonaliser l’opérateur de masse correspondant.
\end{proof}

Si \(F_\ell,F_\nu:\mathbb C^3\to\mathcal H\) sont des repères orthonormaux
complets du même sous-espace physique, alors
\(U_{\mathrm{PMNS}}=F_\ell^*F_\nu\) est unitaire. Si leurs images sont
distinctes, le produit n’est qu’une contraction ; l’unitarité exige
l’égalité des sous-espaces transportés. En présence de dégénérescences,
le mélange est déterminé comme classe double par les changements de repère dans
les blocs dégénérés, non comme une matrice numérique choisie par convention.

La construction minimale de phase
\[
 \Phi_K=I+(e^{i\pi/6}-1)|b_K\rangle\langle b_K|,
 \qquad\|b_K\|=1,
\]
reste un contrôle négatif documenté : elle est unitaire, mais la meilleure
permutation de son identification angulaire donne
\(\theta_{13}\simeq17.02^\circ\), incompatible avec la comparaison de ce
modèle réduit. Le résultat concerne la phase de rang un. Il ne
remplace pas le noyau complet \eqref{vi:vi:eq:gram-neutrinos} et ne permet pas de
présenter comme validée une matrice de mélange qui n’a pas passé son propre
contrôle de rang et d’angles.

\subsubsection{Autoconjugaison neutre et possibilité d’un secteur stérile}

L’alternative Dirac--Majorana utilise la représentation fermionique. Une
structure réelle neutre doit satisfaire
\[
 \mathcal C_\nu^2=I,\qquad
 \mathcal C_\nu M_\nu=M_\nu\mathcal C_\nu,\qquad
 \mathcal C_\nu H_\nu=H_\nu\mathcal C_\nu,
\]
avec \(\mathcal C_\nu\) antiunitaire et une réalification compatible.
La commutation avec la masse préserve sa partie réelle. Pour l’hamiltonien
autoadjoint, en revanche, l’antiunitarité conjugue également le facteur
\(i\) : si \(U_\nu(t)=\exp(-itH_\nu/\hbar)\), alors
\[
 \mathcal C_\nu U_\nu(t)\mathcal C_\nu^{-1}=U_\nu(-t).
\]
Par conséquent, les modes fixes \(\psi=\mathcal C_\nu\psi\) définissent une
structure réelle neutre, mais les conditions écrites ne suffisent pas à
y restreindre toute l’évolution de Schr\"odinger. La conservation de
cette partie fixe requiert que \(\mathcal C_\nu\) commute avec le groupe
\(U_\nu(t)\) ; la proposition~\ref{vi:vi:prop:conjugacion-evolucion}
démontre la condition correspondante sur son générateur. Une
réalisation de Majorana doit déclarer en outre la compatibilité de
l’autoconjugaison avec la représentation fermionique et sa dynamique.
Une charge conservée non nulle qui change
de signe par \(\mathcal C_\nu\) maintient les modes conjugués séparés
dans la réalisation de Dirac. L’involution de mot de la
section~\ref{vi:vi:sec:conjugacion} et la neutralité électrique, considérées
isolément, ne décident pas de cette alternative.

Un secteur stérile requiert un projecteur non nul \(P_s\) tel que
\(P_sP_{\mathrm{act}}=0\) et \(P_sJ_{\text{faible}}=0\).
Si \([P_s,M_\nu]=0\), il est réducteur et ne se mélange pas au secteur actif.
Si \(P_sM_\nu P_{\mathrm{act}}\ne0\), il existe un mélange et le projecteur
cesse de commuter avec la masse. Par conséquent, une profondeur supplémentaire de
l’atlas ne compte pas automatiquement comme quatrième espèce neutrinique.

\subsection{Quarks : saveur, orbite chromatique et transport d’échelle}
\label{vi:vi:subsec:quarks-familias}

L’état quark conserve le n-uplet
\[
 (\sigma_{ud},g,[\gamma]_{\mathrm{col}},\mathfrak M,\mathscr L),
\]
qui distingue le type haut ou bas, la génération, l’orbite de couleur,
la multisection et le registre. La section~\ref{vi:vi:sec:atlas} construit le
résidu chromatique \(\kappa_{\mathrm{col}}(r,c)=r-c\pmod3\) et les
trois classes de \(12\) cellules. L’indice de couleur agit à l’intérieur du
support ; la saveur conserve l’information de génération et de signature que ce
résidu n’identifie pas.

Les six paires angulaires nominales sont
\[
\begin{array}{c|r r|r r}
 f&n_A&s_C&W_+&W_-\\\hline
 u&11&7&73&59\\
 d&17&8&110&94\\
 c&61&8&374&358\\
 s&41&-3&243&249\\
 t&100&-1&599&601\\
 b&71&-5&421&431
\end{array}
\qquad W_\pm=6n_A\pm s_C.
\]
Ces derniers entiers réécrivent le même exposant :
\[
 n_Ax+s_Cy/6
 =\frac{W_+(x+y)+W_-(x-y)}{12}.
\]
Additionner et soustraire \(W_\pm\) permet de retrouver \(n_A\) et \(s_C\), de sorte que
cette présentation conserve les deux canaux. Les six paires appartiennent
à la carte angulaire documentée ; les trois coordonnées restantes
ne sont pas remplies de zéros pour fabriquer une signature complète. Leurs opérateurs
de couleur et les évaluations angulaires sont préservés comme des résultats de
portées différentes.

\begin{proposition}[Dégénérescence chromatique d’une masse de saveur]
\label{vi:vi:prop:masa-color}
Soient \(V_{\mathrm{col}}\simeq\mathbb C^3\) la représentation chromatique
fondamentale et \(U\) son action unitaire irréductible. La réduction
\[
 \mathbb E_{\mathrm{col}}(M)=
 \int_{SU(3)}U(g)MU(g)^*\,dg
\]
est un multiple de \(I_3\). Si \(M\) commute déjà avec l’action, ses trois
composantes chromatiques ont la même valeur de masse de saveur.
\end{proposition}
\begin{proof}
L’invariance de la mesure de Haar donne
\(U(h)\mathbb E_{\mathrm{col}}(M)U(h)^*=
\mathbb E_{\mathrm{col}}(M)\). Le lemme de Schur rend cet
opérateur scalaire. Comme la trace est conservée, le scalaire est
\(\operatorname{Tr}M/3\). Si \(M\) est invariant dès le début,
l’intégrale renvoie \(M\).
\end{proof}

La preuve explique ce qu’élimine le quotient de couleur : la dépendance du
chiffre par rapport au représentant rouge, vert ou bleu. Elle n’élimine pas les
routes chromatiques qui interviennent ensuite dans un état lié. Elle ne
transforme pas non plus cette réduction en sélecteur de saveur.

La masse courante se présente comme une section dans une échelle et un schéma
déclarés, \(m_q(\mu,\mathscr R)\). Si \(t=\log\mu\) et si la connexion
induit \(\gamma_m(g(t))\), le transport satisfait
\[
 \left(\frac{d}{dt}+\gamma_m(g(t))\right)m_q=0,
 \qquad
 m_q(t_2)=\mathcal P\exp\left(-\int_{t_1}^{t_2}
 \gamma_m(g(t))\,dt\right)m_q(t_1).
\]
La connexion et son endomorphisme sont des données du transport, non des nombres
choisis pour égaler un tableau. Un résidu physique n’a de sens
qu’après avoir situé les deux sections dans la même échelle et le même schéma.

Le mélange CKM agit sur une autre coordonnée : le repère de saveur. Une fois
construite la matrice unitaire interne, l’opérateur bas, exprimé dans
le repère haut, est \(B_d^{(u)}=V_{\mathrm{CKM}}B_dV_{\mathrm{CKM}}^*\).
La conjugaison unitaire conserve le polynôme caractéristique, car
\[
 \det(zI-VB_dV^*)=\det\bigl(V(zI-B_d)V^*\bigr)=\det(zI-B_d).
\]
On comprend ainsi la fonction des angles : ils relient des repères, sans
devenir des générateurs rétrospectifs des valeurs propres de masse.

\subsection{Cartes de masse nulle : photon et représentation adjointe de couleur}
\label{vi:vi:subsec:nulo-familias}

Le domaine nul est séparé avant d’appliquer le caractère positif. Dans la
réalisation totale,
\[
 \mathcal H_{\mathrm{phys}}=\mathcal H_0\oplus\mathcal H_+,
 \qquad M|_{\mathcal H_0}=0,\qquad M|_{\mathcal H_+}>0.
\]
L'exponentielle d'une signature finie ne s'annule pas. Une carte nulle
ne s'obtient donc pas en attribuant à une route positive une signature arbitrairement
grande et en appelant zéro son approximation.

Le photon appartient à la carte électromagnétique, avec sa représentation
de jauge et ses deux polarisations transversales physiques. Les gluons
appartiennent à la représentation adjointe chromatique, de dimension \(8\).
L'égalité de leurs masses nulles n'identifie pas les deux représentations.
La fibre, la jauge, les polarisations et les canaux de couplage
continuent de faire partie de l'état.

Une propagation orientée de masse nulle est également compatible. Dans la
carte documentée,
\[
 c_\pm=\frac{c}{1\pm\varepsilon},\qquad
 \varepsilon=q_-^{90}-q_+^{90},\qquad0<\varepsilon<1.
\]
Alors \(c_\pm>0\) et
\[
 \frac{2}{c_+^{-1}+c_-^{-1}}=c.
\]
L'identité découle de l'addition de \((1+\varepsilon)/c\) et
\((1-\varepsilon)/c\). L'asymétrie des deux lectures conserve son
information d'orientation ; elle ne publie pas par elle-même un terme de masse.

\subsection{Bosons massifs et secteur scalaire}
\label{vi:vi:subsec:bosones-familias}

Les descripteurs de \(W\), de \(Z\) et du secteur scalaire ne sont pas contraints
d’entrer dans l’une des \(13\) multisections fermioniques. Ils conservent leur
représentation et leur domaine de réalisation. Les lectures nominales
documentées sont
\[
\begin{array}{c|c|c|c|c}
 X&\text{route nominale}&\text{direction}&(n_A,s_C)&(W_+,W_-)\\\hline
 W&15/37&E\text{--}O&(94,-1)&(563,565)\\
 Z&20/23&E\text{--}O&(95,-1)&(569,571)\\
 H&24/29&N\text{--}S&(97,9)&(591,573)
\end{array}
\]
Ces noms de routes appartiennent au registre nominal conservé. Leur
union avec le domaine orienté complet doit retenir les données de
trajectoire, de représentation et de projecteur. Les décomptes cardinaux
publiés, respectivement \((23,44,11,30)\), \((33,45,10,20)\) et
\((40,34,22,12)\), sont conservés avec l'adresse et ne sont pas remplacés
par une seule étiquette.

Pour \(W\) et \(Z\), on récupère les signatures complètes
\[
 \sigma_W=(94,-1,0,-2,4),\qquad
 \sigma_Z=(95,-1,-11,0,-4).
\]
La lecture \eqref{vi:vi:eq:etapas-e3-sectoriales} produit, pour ces signatures,
\[
 m_W^{[3]}=80378.964909704547024865\ldots\ \mathrm{MeV}/c^2,
\]
\[
 m_Z^{[3]}=91187.533444313407844855\ldots\ \mathrm{MeV}/c^2.
\]
Leur rapport conserve exactement la différence des registres évalués :
\[
 \log\frac{m_Z^{[3]}}{m_W^{[3]}}
 =x-11\Delta_4+2s_{120}-8\zeta_{270}.
\]
En revanche, la coordonnée nominale \((97,9)\) du secteur scalaire appartient
à la carte angulaire \(\mathbb Z^2\) documentée : cette écriture ne lui attribue
ni signature \(\mathbb Z^5\) ni raffinement d'ordre
\(3\). Sa réalisation comme excitation scalaire exige le projecteur
scalaire et les couplages de jauge et fermioniques correspondants.

Pour une espèce instable, le résultat physique est un pôle de l’opérateur
des canaux, non uniquement l’évaluation d’un caractère positif. La
sous-section~\ref{vi:vi:subsec:resonancias-familias} construit cette lecture et
sépare masse de pôle et largeur. Dans le registre documentaire conservé,
les chiffres précédents sont des évaluations conditionnées par la signature ;
l’individualisation de la route, du projecteur et du pôle constitue la réalisation
sectorielle qui doit les accompagner.

\subsection{Grammaire finie des compositions matérielles}
\label{vi:vi:subsec:gramatica-compuestos}

Soient \(\mathcal Q_a\), \(a\in\mathbb Z/3\mathbb Z\), les trois
classes de couleur de l’atlas, chacune de cardinal \(12\). Le dual formel
\(q^\vee\) conserve le support conjugué et change l’orientation de
charge et de couleur. Ce n’est pas la simple cellule \(\rho_{\mathrm c}(q)\) dépourvue
de ces données. Les expressions mésoniques et baryoniques sont définies par
\begin{align}
 \mathfrak M&=\{q\otimes(q')^\vee:
 q,q'\in\mathcal Q_a\text{ pour un certain }a\},\notag\\
 \mathfrak B&=\mathcal Q_0\times\mathcal Q_1\times\mathcal Q_2.
 \label{vi:vi:eq:gramatica-meson-barion}
\end{align}
Les paires sont ordonnées et les positions de couleur du triplet sont
fixées. La neutralité résiduelle vaut, respectivement,
\(a-a=0\) et \(0+1+2=0\pmod3\).

\begin{proposition}[Recensement compositionnel interne]
\label{vi:vi:prop:censo-compuestos}
La grammaire contient \(432\) expressions mésoniques, \(1728\)
expressions baryoniques et \(372\) flèches chargées compatibles, réparties
en \(320\) leptoniques et \(52\) quark. Ces cardinaux comptent
des expressions et des flèches internes, non des espèces physiques ni des prédictions
indépendantes de masses.
\end{proposition}
\begin{proof}
Pour chaque couleur, on choisit une paire ordonnée de \(12\) cellules ; par conséquent,
\(|\mathfrak M|=3\cdot12^2=432\). Les trois positions du baryon donnent
\(|\mathfrak B|=12^3=1728\). Les transitions lepton--neutrino de
même profondeur ne partagent que \(G_2,G_4\) ; dans les deux directions, on
obtient \(2(8\cdot4+16\cdot8)=320\).

Pour les transitions haut--bas, on exige en outre l’égalité de couleur.
Dans \(G_1\), les multiplicités chromatiques sont \((2,1,1)\) et
\((0,1,1)\) ; dans \(G_3\), ce sont \((4,4,4)\) et \((2,2,2)\).
Il y a \(2+24=26\) paires d’une direction et \(52\) des deux. Ajouter
\(320+52\) complète le recensement. Tous les facteurs proviennent de l’atlas
antérieur à l’évaluation des masses.
\end{proof}

Les \(81\) identités des cellules complètent les flèches neutres
de cette grammaire, mais une identité compositionnelle ne devient pas pour
autant un boson neutre. La représentation et le canal de réalisation
sont les données qui distinguent les deux notions.

\subsection{Composition du registre et mémoire de liaison}
\label{vi:vi:subsec:composicion-registro}

Un composé contient plus d’informations que la liste de ses
constituants. Soient \(L_1,L_2\) des registres enrichis, \(g\) le
transport qui aligne phase, feuillet et orientation, et \(B\) la sous-trace
identifiée à la frontière commune. Dans les coordonnées d’action et
d’événements, la composition est
\begin{equation}
 (n,e)_1\star_{(g,B)}(n,e)_2
 =\bigl(n_1+n_2-n_B,\ e_1+ge_2-e_B\bigr).
 \label{vi:vi:eq:composicion-frontera}
\end{equation}
La soustraction évite de compter deux fois une même sous-trace ; le transport
conserve l’orientation à l’interface. Les prédécesseurs torsionnels et
les marques de couronne se composent au même niveau avant d’appliquer de nouveau
\(L_{\mathrm{tick}}\). La décomposition réversible des événements
de la proposition~\ref{vi:vi:prop:eventos} permet de transporter le registre
complet sans le réduire à son total.

Pour un arbre ordonné \(\mathcal T\) de \(n\) constituants, on obtient
\begin{equation}
 (\Gamma_1,\ldots,\Gamma_n;\mathcal T,\partial)
 \xmapsto{\mathfrak C_{12}^{(n)}}L_Y
 \xmapsto{L_{\mathrm{tick}}}\sigma_Y
 \longmapsto (\eta_Y,M_Y).
 \label{vi:vi:eq:cadena-compuesto}
\end{equation}
La frontière \(\partial\) comprend les transports et les sous-traces de chaque
arête de l’arbre. Une fois fixées ces données admissibles, chaque composition
est une opération explicite. Le caractère s’applique à la signature finale ;
il ne sélectionne pas préalablement l’arbre par la masse d’arrivée.

\begin{proposition}[Compatibilité du parenthésage]
\label{vi:vi:prop:cociclo-compuesto}
Dans une carte de registres où
\(L_1\star L_2=L_1+L_2+\mathfrak b(L_1,L_2)\), les deux
parenthésages d'un triplet admissible coïncident si et seulement si
\begin{align*}
 \mathfrak b(L_1,L_2)+\mathfrak b(L_1\star L_2,L_3)
 =\mathfrak b(L_2,L_3)+\mathfrak b(L_1,L_2\star L_3).
\end{align*}
L’égalité est exigée dans le domaine où les deux arbres représentent la
même interface orientée.
\end{proposition}
\begin{proof}
Développer \((L_1\star L_2)\star L_3\) donne
\(L_1+L_2+L_3+\mathfrak b(L_1,L_2)+\mathfrak b(L_1\star L_2,L_3)\).
Développer \(L_1\star(L_2\star L_3)\) donne la même somme de trois
registres et les deux termes du second membre. Leur égalité équivaut
exactement à l’identité énoncée.
\end{proof}

La loi cocyclique assure l’indépendance d’un parenthésage équivalent ;
elle n’identifie pas des arbres physiquement différents. Deux structures de liaison
qui conservent des canaux différents restent des réalisations distinctes
jusqu’à disposer d’un opérateur d’entrelacement des registres, des projecteurs et des opérateurs.

\subsubsection{Distinction entre somme des signatures et registre intégral de composition}

L’extracteur réduit historique d’événements se définit par
\[
 F_0(n,e)=\left(n,\sum_m e_m,K_0(e),
 \sum_{a=0}^3\operatorname{sgn}(e_a+e_{a+4}+e_{a+8}),
 \sum_m\tau_m\tau_{m+3}\right),
\]
où \(\tau_m=\operatorname{sgn}(e_m)\) et
\(K_0(e)=\sum_{\tau_m=0}\Sigma_{12}(m)\), avec
\(\Sigma_{12}=(+,+,+,-,-,-)^2\). Sa troisième coordonnée est une
lecture d’événements nuls ; ce n’est pas la torsion forte
\(k_{\Delta_4}=T_z-2C_z\) de la
section~\ref{vi:vi:sec:operador-masa}.

Considérons les registres des routes
\[
 a=(1,3,-1,+1),\quad b=(6,7,-1,+1),\quad c=(1,2,-1,-1),
\]
avec
\begin{align*}
 e(a)&=(5,4,4,-4,-4,-6)^2,\\
 e(b)&=(4,5,4,-4,-6,-4)^2,\\
 e(c)&=(1,4,1,-2,-4,-3)^2.
\end{align*}
Ils ont \(F_0(a)=F_0(b)=(8,-2,0,0,-12)\) et
\(F_0(c)=(28,-6,0,-4,-12)\). En superposant les registres,
\[
 F_0(a\boxplus c)=(36,-8,0,0,-12),\qquad
 F_0(b\boxplus c)=(36,-8,0,-2,-12).
\]
Les entrées projetées sont égales et les sorties composées diffèrent.
Par conséquent, aucune opération sur les deux quintuplets de cette
projection ne récupère toutes les superpositions. L’exemple démontre une
obstruction précise pour \(F_0\) ; il n’attribue pas le même résultat à toute
extension possible d’une signature. La construction en vigueur évite cette
perte au moyen de \eqref{vi:vi:eq:cadena-compuesto} : elle conserve les événements
et les prédécesseurs, compose et alors seulement évalue l’extracteur complet.

\subsection{Mésons : paire duale, singulet et différenciation des excitations}
\label{vi:vi:subsec:mesones-familias}

L’expression \(q\otimes(q')^\vee\) établit la paire ordonnée et sa
neutralité résiduelle. Le singulet de couleur appartient à la réalisation
\(V\otimes\overline V\), \(V=\mathbb C^3\). Si
\(\Omega=\sum_{a=1}^3|a\bar a\rangle\), le projecteur est
\begin{equation}
 P_M^{(1)}=\frac13|\Omega\rangle\langle\Omega|.
 \label{vi:vi:eq:singlete-meson}
\end{equation}

\begin{proposition}[Projection mésonique de couleur]
L’application \(P_M^{(1)}\) est un projecteur orthogonal de rang un et
son image est invariante sous \(U\otimes\overline U\), \(U\in SU(3)\).
\end{proposition}
\begin{proof}
On a \(\langle\Omega,\Omega\rangle=3\), donc
\((P_M^{(1)})^2=P_M^{(1)}\). De plus,
\[
 (U\otimes\overline U)\Omega
 =\sum_{a,b,c}U_{ba}\overline U_{ca}|b\bar c\rangle
 =\sum_b|b\bar b\rangle=\Omega.
\]
L’image conserve ainsi la contraction chromatique sans privilégier une couleur.
\end{proof}

Le descripteur mésonique réunit saveurs, spin et parité, isospin et charge,
excitation radiale ou orbitale, arbre de liaison et frontière. L’étiquette
\(J^{PC}\) est utilisée lorsque \(C\) est défini pour le canal ; on
n’attribue pas automatiquement de conjugaison de charge à un état chargé. Deux
états ayant les mêmes saveurs peuvent avoir des projecteurs
radiaux, orbitaux ou de canal différents et, par conséquent, des spectres différents. La
construction orbitale de la section~\ref{vi:vi:sec:kepler} apporte une
lecture possible lorsque ses hypothèses dynamiques sont satisfaites ; les nombres
orbitaux ne se déduisent pas du seul nom du méson.

Le registre des évaluations composées conserve
\[
\begin{array}{c|r r r r r}
 &n_A&s_C&k_{\Delta_4}&b_{120}&b_\zeta\\\hline
 \pi^0&44&-4&-25&1&-2\\
 \pi^\pm&44&1&-2&2&-1\\
 K^\pm&54&-1&17&-2&2\\
 K^0&54&0&32&3&-2
\end{array}
\]
Pour ces signatures, la même lecture \eqref{vi:vi:eq:etapas-e3-sectoriales}
explique les différences des paires :
\begin{align}
 \log\frac{m_{\pi^\pm}^{[3]}}{m_{\pi^0}^{[3]}}
 &=\frac56y+23\Delta_4+s_{120}+\zeta_{270},\notag\\
 \log\frac{m_{K^0}^{[3]}}{m_{K^\pm}^{[3]}}
 &=\frac16y+15\Delta_4+5s_{120}-4\zeta_{270}.
 \label{vi:vi:eq:desdoblamientos-mesonicos}
\end{align}
Les termes d’action \(44x\) et \(54x\) s’annulent au sein de
chaque paire. Les différences restantes conservent quelle coordonnée de
registre produit chaque contribution. Ce sont des conséquences exactes des
signatures déclarées, non une preuve de substitution de leur génération par l’arbre
de liaison.

Les évaluations résultantes sont
\begin{align*}
 m_{\pi^0}^{[3]}&=134.976724327872125739\ldots\ \mathrm{MeV}/c^2,\\
 m_{\pi^\pm}^{[3]}&=139.570485171572191374\ldots\ \mathrm{MeV}/c^2,\\
 m_{K^\pm}^{[3]}&=493.677085649530612475\ldots\ \mathrm{MeV}/c^2,\\
 m_{K^0}^{[3]}&=497.611186497750457358\ldots\ \mathrm{MeV}/c^2.
\end{align*}
Chaque chiffre conserve le
statut d’évaluation exacte conditionnée par la signature composée
récupérée. Le singulet, la grammaire et l’extraction du registre sont
construits ; la réalisation nominale complète exige que l’arbre de
liaison régénère cette signature avant d’ouvrir la comparaison.

\subsection{Baryons : triplet chromatique, échange et différence neutron--proton}
\label{vi:vi:subsec:bariones-familias}

La neutralité de \(q_0\otimes q_1\otimes q_2\) conserve les trois
positions chromatiques. Dans la réalisation tensorielle, on introduit le vecteur
normalisé
\[
 |\varepsilon\rangle=\frac1{\sqrt6}
 \sum_{a,b,c=1}^3\varepsilon_{abc}|abc\rangle,
 \qquad P_B^{(1)}=|\varepsilon\rangle\langle\varepsilon|.
\]

\begin{proposition}[Singulet baryonique et règle d’échange]
\label{vi:vi:prop:singlete-barion}
L’opérateur \(P_B^{(1)}\) est un projecteur orthogonal de rang un,
invariant sous \(U^{\otimes3}\), \(U\in SU(3)\). Une permutation
de deux positions change le signe de son vecteur générateur et conserve
le projecteur.
\end{proposition}
\begin{proof}
Il y a \(6\) termes non nuls orthogonaux, chacun de module
\(1/\sqrt6\) ; le vecteur a donc une norme égale à un. L’identité
du déterminant donne
\(U^{\otimes3}|\varepsilon\rangle=
\det(U)|\varepsilon\rangle=|\varepsilon\rangle\).
Une transposition change le signe de \(\varepsilon_{abc}\). Dans le
produit extérieur du vecteur par son adjoint, les deux signes s’annulent.
\end{proof}

L’antisymétrie chromatique est une partie de l’état complet. Le projecteur
physique réunit couleur, spin et parité, isospin, charge et symétrie
d’échange des composantes orbitales et de saveur. Le principe de
Pauli agit dans le produit fermionique complet ; il ne suffit pas de l’appliquer à une
simple liste de couleurs. L’arbre triple enregistre quelle paire a été couplée
en premier et comment le troisième constituant a traversé la frontière. La
proposition~\ref{vi:vi:prop:cociclo-compuesto} contrôle les
parenthésages équivalents.

Les signatures récupérées du proton et du neutron sont
\[
 \sigma_p=(59,0,9,1,0),\qquad
 \sigma_n=(59,1,-34,0,1).
\]
Ils partagent la coordonnée d’action \(59\), mais diffèrent en feuillet,
torsion et retours. Leur différence d’exposant est
\begin{equation}
 \log\frac{m_n^{[3]}}{m_p^{[3]}}
 =\frac16y-43\Delta_4-s_{120}+\zeta_{270}.
 \label{vi:vi:eq:diferencia-neutron-proton}
\end{equation}
L’égalité se démontre en soustrayant les cinq coordonnées. Dans cette signature,
\(-34\) est un entier torsionnel du neutron ; il ne s’identifie pas à
l’exposant décimal \(-34\) de l’échelle absolue d’action. L’égalité
de deux entiers dans des applications différentes n’élimine ni leur type ni leur provenance.

Les évaluations conservées sont
\begin{align*}
 m_p^{[3]}&=938.271751546984898298\ldots\ \mathrm{MeV}/c^2,\\
 m_n^{[3]}&=939.565810472507023312\ldots\ \mathrm{MeV}/c^2.
\end{align*}
La formule \eqref{vi:vi:eq:diferencia-neutron-proton} localise l'effet de
chaque composante des signatures nominales. La dérivation physique de la
différence exige de transporter l'isospin, l'orientation des lacunes,
la frontière de liaison et la contribution électromagnétique dans une même carte.
L'évaluation du caractère conserve son exactitude sous la signature fixée ;
cette exactitude ne remplace pas la construction prospective de la liaison.

\subsection{Compositions d’ordre supérieur et équivalence des canaux}
\label{vi:vi:subsec:compuestos-superiores}

L’application \(\mathfrak C_{12}^{(n)}\) admet un nombre arbitraire
fini de constituants. Un tétraquark emploie \(qq\bar q\bar q\) ;
un pentaquark, \(qqqq\bar q\) ; un état gluonique lié utilise
des représentations adjointes projetées sur le singulet ; un hybride conserve
des constituants fermioniques et de jauge dans le même arbre. Noyaux,
atomes et molécules emploient la même antériorité du registre sur la lecture,
avec leurs propres couplages et frontières.

L’énergie de liaison appartient à la réalisation du registre composé.
Elle ne s’obtient pas en additionnant des masses isolées et en baptisant liaison le
résidu souhaité. Même lorsqu’un cocycle prend des valeurs additives sur
les registres, l’extracteur de signes et de corrélations s’applique ensuite ;
c’est pourquoi le caractère multiplicatif d’une signature ne transforme pas la masse
composée en somme de masses constituantes.

Si deux arbres \(\mathcal T_1,\mathcal T_2\) possèdent un opérateur unitaire
\(V\) qui entrelace leurs opérateurs et leurs projecteurs de canal,
\[
 VH_1=H_2V,\qquad VP_1=P_2V,
\]
alors \(V(z-H_1)^{-1}=(z-H_2)^{-1}V\) dans l'ensemble résolvant commun.
Son prolongement, lorsqu'il existe et est transporté dans la même carte,
conserve les pôles et leurs multiplicités. Cette égalité donne le critère
opérationnel pour identifier deux réalisations d'un composite ; partager
des saveurs ou une masse approximative ne suffit pas.

\subsection{Résonances : masse de pôle, largeur et mémoire de survie}
\label{vi:vi:subsec:resonancias-familias}

Soient \(P\) le sous-espace préparé et \(Q=I-P\) celui des canaux
éliminés. Dans une région où \(z-QHQ\) est inversible, l’élimination
de la composante \(Q\psi\) de \((z-H)\psi=0\) donne
\begin{equation}
 H_{\mathrm{eff}}(z)=PHP+PHQ(z-QHQ)^{-1}QHP.
 \label{vi:vi:eq:feshbach-familias}
\end{equation}

\begin{proposition}[Lecture réduite d’un état lié ou résonant]
\label{vi:vi:prop:feshbach-familias}
Si \(P\mathcal H\) est de dimension finie, les valeurs spectrales
admissibles du problème réduit satisfont
\(\det(z-H_{\mathrm{eff}}(z))=0\). Une résonance exige, en outre,
un prolongement méromorphe à la carte de canaux pertinente et un zéro
isolé de résidu physique admissible.
\end{proposition}
\begin{proof}
L’équation dans le complément est
\((z-QHQ)Q\psi=QHP\psi\). La résoudre et substituer dans l’équation
de \(P\psi\) produit \((z-H_{\mathrm{eff}}(z))P\psi=0\).
En dimension finie, il existe une solution non nulle exactement lorsque
le déterminant s’annule. Le calcul algébrique se réalise dans le domaine
de la résolvante ; le passage à une résonance utilise le prolongement et le
résidu indiqués, qui ne se déduisent pas de la seule ouverture d’un canal.
\end{proof}

Un état stable est une raie spectrale normalisable sans canal de
désintégration ouvert ; un état lié est sous le seuil
pertinent, même s’il peut se désintégrer par un autre canal ; une résonance est un
pôle du prolongement. Si
\[
 z_X=M_Xc_{\mathrm{int}}^2-\frac{i}{2}\Gamma_X,
 \qquad\Gamma_X>0,
\]
la partie réelle et la partie imaginaire sont deux lecteurs coordonnés du même
pôle. La largeur n’est pas une sixième coordonnée de \(\sigma\).

La lecture discrète utilise une valeur propre de survie
\(\lambda_X=re^{-i\theta}\), \(0<r<1\), et un pas temporel
\(t_0>0\). La branche \(\theta\in I_\theta\) de la carte de
quasi-énergie étant fixée,
\begin{equation}
 z_X=\frac{\hbar_{\mathrm{int}}}{t_0}
 (\theta+i\log r),\quad
 M_X=\frac{\hbar_{\mathrm{int}}\theta}{t_0c_{\mathrm{int}}^2},\quad
 \Gamma_X=-\frac{2\hbar_{\mathrm{int}}}{t_0}\log r.
 \label{vi:vi:eq:anchura-discreta-familias}
\end{equation}
En effet, substituer dans
\(\lambda_X=\exp(-iz_Xt_0/\hbar_{\mathrm{int}})\) récupère
\(re^{-i\theta}\). Comme \(\log r<0\), le pôle est retardé et la
largeur positive. La probabilité de survie après \(N\) pas est
\(r^{2N}=\exp(-\Gamma_XNt_0/\hbar_{\mathrm{int}})\), d’où
\(\tau_X=\hbar_{\mathrm{int}}/\Gamma_X\).

Si cette lecture est obtenue à partir d'une compression \(SUS\), l'interprétation
asymptotique par un mode unique exige une valeur propre dominante
simple, un recouvrement initial non nul et une séparation en module du reste
du spectre. Une normalisation sur les probabilités, au lieu des
amplitudes, modifie le facteur \(2\). La branche temporelle et cette convention
accompagnent donc chaque largeur publiée.

L’égalité de masse et de largeur d’états conjugués requiert de transporter
également les canaux :
\[
 U_C H_{\mathrm{eff}}^X(z)U_C^{-1}=H_{\mathrm{eff}}^{\bar X}(z).
\]
Alors, les mêmes pôles sont transportés avec la même prescription
causale. Une représentation antiunitaire exige de transporter en outre cette
prescription : la conjugaison seule de \(z\) échange les pôles
retardés et avancés.

\subsection{Trois codomaines topologiques et secteur virtuel}
\label{vi:vi:subsec:topologicos-familias}

Les classes topologiques conservent des objets et des règles de fusion, non une
masse scalaire universelle. La donnée primaire est
\((N_{ab}^{c},F^{abc}_d,R^{ab}_c)\), avec ses équations de cohérence.
Un écart d’énergie requiert ensuite une réalisation hamiltonienne et
un projecteur spectral. L’attribution de ce type de sortie permet
d’inclure ces secteurs sans les transformer en lignes de masse élémentaire.

\subsubsection{Sector Fibonacci}

Les objets \(\mathbf1,\tau\) satisfont
\(\tau\otimes\tau=\mathbf1\oplus\tau\). La dimension positive
\(d_\tau\) conserve la multiplication de fusion :
\(d_\tau^2=1+d_\tau\). La réalisation du corpus identifie cette
dimension à la coordonnée interne \(\varphi\) déjà engendrée ;
l’équation de fusion est sa reconnaissance dans ce codomaine, non un
substitut de sa généalogie APP--TRIT--TPK. Dans la convention fixée,
\[
 F_\tau^{\tau\tau\tau}=\begin{pmatrix}
 \varphi^{-1}&\varphi^{-1/2}\\
 \varphi^{-1/2}&-\varphi^{-1}
 \end{pmatrix},\qquad
 R_{\mathbf1}^{\tau\tau}=e^{-4\pi i/5},\quad
 R_\tau^{\tau\tau}=e^{3\pi i/5}.
\]
L’identité \(\varphi^{-2}+\varphi^{-1}=1\) prouve directement
que \(F^*F=I\) ; les termes non diagonaux s’annulent. Les phases
de \(R\) conservent l’orientation du tressage.

\subsubsection{Secteur Ising et réalisation de Majorana}

Les règles sont
\[
 \psi\otimes\psi=\mathbf1,\quad
 \psi\otimes\sigma=\sigma,\quad
 \sigma\otimes\sigma=\mathbf1\oplus\psi,
\]
avec
\[
 F_\sigma^{\sigma\sigma\sigma}=\frac1{\sqrt2}
 \begin{pmatrix}1&1\\1&-1\end{pmatrix},\qquad
 R_{\mathbf1}^{\sigma\sigma}=e^{-\pi i/8},\quad
 R_\psi^{\sigma\sigma}=e^{3\pi i/8}.
\]
La multiplication matricielle donne \(F^2=I\). Le secteur conserve ses
trois types de fusion et sa représentation ; sa dénomination Majorana
ne transforme pas automatiquement un neutrino du secteur précédent en mode
autoconjugué. Cette affirmation utilise les conditions antiunitaires
spécifiques de la réalisation neutre.

\subsubsection{Secteur abélien de six classes}

Pour \(a,b\in\mathbb Z/6\mathbb Z\),
\[
 a\otimes b=a+b\pmod6,\qquad
 R_{ab}=\zeta_6^{ab},\qquad\zeta_6=e^{2\pi i/6}.
\]
L’identité \(R_{a+b,c}=R_{a,c}R_{b,c}\) et son analogue dans le
second argument expriment le caractère bilinéaire du tressage.
La construction dupliquée
\(\mathcal C\boxtimes\mathcal C^{\mathrm{rev}}\) conserve les deux
orientations. Ses classes ne doivent pas être comptées comme six nouvelles masses
élémentaires.

\subsubsection{Persistance finie et virtualité}

Le secteur virtuel conserve des extensions finies, leur canal et leur horizon
de persistance. Une section finie avec obstruction terminale n’occupe pas
le même codomaine qu’une section asymptotique persistante. La
virtualité se lit dans une amplitude ou une résolvante et n’équivaut pas, par
définition, à une masse imaginaire universelle. Il ne suffit pas non plus qu’une
extension globale ne soit pas certifiée dans une fiche pour affirmer qu’il
existe une obstruction : l’horizon et sa preuve sont des données différentes.
La classification conserve précisément cette information finie sans
inventer une particule stable pour chaque ligne interne de calcul.

\subsection{Classification en vingt-trois classes et portée de leurs réalisations}
\label{vi:vi:subsec:23-clases}

L'inventaire public réunit \(17\) classes élémentaires, \(2\) familles
composites, \(3\) codomaines topologiques et \(1\) secteur virtuel :
\(17+2+3+1=23\). Les trois lignes neutriniques nomment des saveurs, les
six lignes de quarks conservent leur triplicité chromatique et la ligne gluonique
conserve la représentation adjointe. Aucune de ces conventions de
décompte ne modifie l'atlas \(81/56/13/324\).

La liste suivante maintient le type de sortie et la portée des
réalisations documentées. Les expressions « signature déclarée » et
« opérateur de fibre » renvoient aux constructions explicites de cette
section, et non à des degrés différents de précision décimale.
Le registre de provenance distingue la fermeture scalaire électronique et
les deux cartes nulles de la scalarisation physique encore en attente dans
les familles non centrales. Dans les neutrinos, il conserve en attente la
réalisation complète du mélange et de l'échelle ; dans les quarks, la sélection de
saveur et sa carte d'échelle ; dans les bosons nominaux, l'union de la route
et du pôle ; dans les composites nommés, la réalisation scalaire de leur
liaison. Ces états documentaires ne modifient ni les opérateurs, ni les
grammaires, ni les évaluations conditionnées qui ont été exposées.

\begin{longtable}{@{}p{.11\linewidth}p{.20\linewidth}p{.62\linewidth}@{}}
 \toprule
 Classe&Représentants&Résultat conservé et opération de réalisation\\
 \midrule\endfirsthead
 \toprule Classe&Représentants&Résultat conservé et opération de réalisation\\
 \midrule\endhead
 \bottomrule\endfoot
 1&\(e^-,e^+\)&Fermeture bêta scalaire de rang un et représentation conjuguée.\\
 2&\(\mu^-,\mu^+\)&Opérateurs exacts de fibre ; évaluation nominale et raffinement d’ordre \(3\) avec signature déclarée ; sélection spectrale physique différenciée.\\
 3&\(\tau^-,\tau^+\)&Opérateurs exacts de fibre ; évaluation nominale et raffinement d’ordre \(3\) avec signature déclarée ; sélection spectrale physique différenciée.\\
 4&\(\nu_e,\bar\nu_e\)&Opérateur neutre de fibre ; réalisation du mélange, du repère de masse et de l’échelle physique par leurs applications complètes.\\
 5&\(\nu_\mu,\bar\nu_\mu\)&Même domaine neutre, avec une étiquette de saveur différente ; non une profondeur individuelle de l’atlas.\\
 6&\(\nu_\tau,\bar\nu_\tau\)&Même domaine neutre ; conservation de l’alternative Dirac--Majorana sous ses conditions.\\
 7&\(u,\bar u\)&Opérateur de fibre chromatique et lecture angulaire nominale ; scalarisation, schéma et échelle conservés.\\
 8&\(d,\bar d\)&Opérateur de fibre chromatique et lecture angulaire nominale ; scalarisation, schéma et échelle conservés.\\
 9&\(c,\bar c\)&Opérateur de fibre chromatique et lecture angulaire nominale ; scalarisation, schéma et échelle conservés.\\
 10&\(s,\bar s\)&Opérateur de fibre chromatique et lecture angulaire nominale ; scalarisation, schéma et échelle conservés.\\
 11&\(t,\bar t\)&Opérateur de fibre chromatique et lecture angulaire nominale ; convention de masse et canal de comparaison explicites.\\
 12&\(b,\bar b\)&Opérateur de fibre chromatique et lecture angulaire nominale ; scalarisation, schéma et échelle conservés.\\
 13&Photon&Section de masse exactement nulle ; représentation électromagnétique.\\
 14&Gluons&Section de masse exactement nulle ; représentation adjointe de couleur.\\
 15&\(W^+,W^-\)&Signature nominale complète et évaluation d’ordre \(3\) ; réunion avec route, projecteur et pôle comme réalisation physique différenciée.\\
 16&\(Z^0\)&Signature nominale complète et évaluation d’ordre \(3\) ; réunion avec route, projecteur et pôle comme réalisation physique différenciée.\\
 17&\(H\)&Descripteur scalaire et évaluation angulaire nominale ; projecteur, couplages et pôle de la réalisation scalaire.\\
 18&Mésons&Grammaire exacte de paire duale et de liaison ; signatures composées récupérées ; scalarisation nominale et canaux de pôle différenciés.\\
 19&Baryons&Grammaire exacte de triplet chromatique et de liaison ; signatures du proton et du neutron récupérées ; réalisation nominale de liaison différenciée.\\
 20&Fibonacci
&Codomaine de fusion et de tressage ; un écart requiert sa réalisation hamiltonienne.\\
 21&Ising/Majorana&Codomaine de fusion et structure réelle selon la représentation ; sans masse scalaire universelle attribuée.\\
 22&\(\mathbb Z/6\mathbb Z\)&Six classes de tressage ; codomaine topologique, non six masses universelles.\\
 23&Secciones virtuales&Extension finie et horizon d’obstruction ; sans scalaire stable universel.\\
\end{longtable}

Les états documentaires exacts sont publiés dans l'annexe~B.3 : le
tableau~\ref{vi:tab:vi-estados-clases} conserve la légende des statuts et le
tableau~\ref{vi:tab:vi-clases} conserve, pour les vingt-trois classes, résidence,
incidence, représentation et état scalaire.

\begin{proposition}[Produit fibré de réalisation et codomaine observable]
\label{vi:vi:prop:producto-fibrado-especie-ruta}
Pour \(X\in\mathcal P_{\rm adm}\), la fibre définie dans
\eqref{vi:vi:eq:fibra-descriptor} est non vide si et seulement si
\[
 d_{\mathcal P}(X)\in\operatorname{im}d_{\mathcal R}.
\]
Sur cette fibre, la mesure spectrale conjointe
\eqref{vi:vi:eq:familias-pvm} détermine le codomaine observable des lecteurs
\((M_X^{\beta,D},M_X^{\beta,\Omega})\). Une préparation \(P_X\) publie une
raie scalaire \(m\) exactement lorsqu'elle satisfait le critère de la
proposition~\ref{vi:vi:prop:familias-linea} ; dans les autres cas, elle conserve le
multiplet, la distribution spectrale ou le type sectoriel produit.
\end{proposition}
\begin{proof}
Par la définition de l’image, la condition écrite équivaut à l’existence
d’une route \(\gamma\in\mathcal R_{\rm adm}\) telle que
\(d_{\mathcal R}(\gamma)=d_{\mathcal P}(X)\) ; c’est précisément la
condition d’appartenance à \(F_X\). La seconde affirmation découle du théorème
spectral appliqué aux deux opérateurs diagonaux et de la caractérisation
du support ponctuel démontrée dans la
proposition~\ref{vi:vi:prop:familias-linea}.
\end{proof}

L'état documentaire externe contient \(471\) observables de masse et
\(384\) de largeur, soit \(855\) au total. Ce sont des lignes de comparaison, et non
\(855\) prédictions internes. En particulier, les \(243\) lignes
mésoniques de masse et les \(166\) lignes baryoniques ne coïncident pas avec les cardinaux
\(432\) et \(1728\) de la grammaire : elles comptent des objets de types distincts.
Le tableau externe est ouvert après avoir fixé identité, fibre, registre,
opérateur, projecteur et carte.

La reconstruction d’un rapport positif au moyen d’une tour qui reçoit
ce rapport démontre la capacité représentative du système d’évaluation.
La sortie indépendante d’une espèce exige, en revanche, que les
coefficients de la tour proviennent de sa route et de sa frontière. Cette
séparation conserve aussi bien la complétude représentative que les
résultats sectoriels effectifs. L’ensemble est organisé par une
règle commune : construire la structure et sa mémoire, appliquer le lecteur
du type correspondant et effectuer ensuite la reconnaissance physique.
